\subsection*{3.5 The Gradient, the Infrastructure Layer, and the Limits of Prediction}

The empirical shift described above is best understood as a change in the economic gradient that powered the installation phase. During installation, the gradient was steep: frontier capability commanded a premium, and capital invested in scaling produced capability gains that could be monetized. That gradient is now flattening from both ends. From below, falling inference costs make commoditized capability sufficient for most deployment tasks. From above, the cost of frontier participation rises faster than the return, as scaling yields diminishing capability gains per unit of compute.

These are two distinct gradients, not one. The first is the gap between frontier capability and commoditized capability. The second is the gap between frontier cost and frontier return. They interact---commoditization from below reduces the return on frontier investment from above---but they are separate variables, and the paper treats them as related pressures rather than as a single measured quantity. Conflating them would make the argument appear to measure one thing while describing another.

A system facing a flattening gradient may respond in several ways: exit, consolidation, renewed investment, technological discontinuity, or a shift toward maintaining the conditions under which the gradient remains viable. The paper does not claim that maintenance is the necessary outcome for every firm or every layer. It claims that maintenance is the rational response for firms positioned in an infrastructure layer whose demand is long-term and not exhausted by commoditization---the layer through which deployment traffic passes, rather than the layer at which capability is differentiated.

The distinction is between product logic and infrastructure logic. Product logic extracts margin from differentiation; when differentiation declines, the rational responses are exit, consolidation, or a new product. Infrastructure logic extracts rent from guaranteed access; when differentiation declines but demand for the layer remains long-term and general, the rational response is to control the conditions of access to the layer. The margin is lower, but it is durable. This is why firms invest in electricity generation at thin margins: not because the product is differentiated, but because the layer is indispensable and access to it is a long-term guarantee. The rent is the return on controlling a layer that others must traverse, not a failure of competition.

Applied to AI, the frontier model is a product; the distribution, compute, and governance layers are infrastructure. Commoditization of the product layer---falling inference costs, open-weight models---does not eliminate the infrastructure layer. It makes the infrastructure layer the site of rent. Maintenance, in this sense, is the rent-seeking behavior of firms that have moved, or are moving, from product competition to infrastructure position---not a mysterious systemic tendency. The conversion can be observed in microcosm. OpenAI's 2019 partnership with Microsoft---exclusive cloud rights, Azure OpenAI Service as a distribution layer, Microsoft's 49\% share of economic returns---is a textbook case of a frontier lab converting a product position (GPT-class models) into an infrastructure position (Azure's enterprise AI distribution). Anthropic's partnerships with Amazon Bedrock and Google Cloud Vertex AI follow the same pattern. The conversion is not concealed; it is disclosed in financial filings. What is not yet documented is whether governance arrangements (evaluator access, certification, federal preemption) are being designed to protect the converted position rather than the original product. That is the question Section~4 takes up.

The paper uses the term \textbf{attractor} heuristically to denote this family of stabilizing responses---arrangements that preserve the conditions under which frontier production remains economically viable. The term is borrowed from dynamical-systems theory, where it denotes a state toward which a system evolves under specified dynamics; here it is used only as a label for the class of stabilizing responses, with no formal dynamical model implied. The paper does not specify the state variables, feedback mechanisms, or conditions under which such an attractor would exist in the formal sense. It offers a vocabulary for describing the pressure and the class of responses, not a model that predicts which response will be selected.

This distinction matters for what the paper can claim. The pressure is observable: falling inference costs, rising frontier costs, migration of deployment to commoditized models. The response is contingent: multiple actors with competing interests, some observable and some not, shape the trajectory. The paper evaluates the design features of one response without claiming that this response was inevitable or that its form was determined. Whether the frontier laboratories are actually converting a product position into an infrastructure position is an empirical question, addressed in Section~4, not a theoretical deduction.
